% ------------------------------------------------------------------------------
% SECCIÓN 5: GLOSARIO Y NOTAS TÉCNICAS
% ------------------------------------------------------------------------------
\section{Glosario y Notas Técnicas}

Para garantizar un entendimiento inequívoco entre los equipos de desarrollo, analistas de negocio y los tomadores de decisión, se definen los siguientes términos técnicos clave:

\begin{GlossaryBox}{CRM (\textit{Customer Relationship Management})}
    Plataforma de software orientada a la gestión integrada de las relaciones comerciales, ventas, seguimiento de prospectos y atención al cliente, cuyo fin es consolidar toda la interacción en una base unificada.
\end{GlossaryBox}

\begin{GlossaryBox}{ERP (\textit{Enterprise Resource Planning})}
    Sistema de planificación de recursos empresariales encargado de centralizar y automatizar procesos contables, facturación, compras, inventarios y logística dentro de una organización.
\end{GlossaryBox}

\begin{GlossaryBox}{API RESTful (\textit{Representational State Transfer})}
    Interfaz de comunicación basada en los estándares del protocolo HTTP (GET, POST, PUT, DELETE), con arquitectura sin estado (\textit{stateless}) y payloads usualmente formateados en JSON, diseñada para interoperabilidad entre sistemas distribuidos.
\end{GlossaryBox}

\begin{GlossaryBox}{OAuth 2.0}
    Protocolo de autorización estándar de la industria que permite a aplicaciones de terceros obtener acceso limitado y seguro a servicios de usuario (por ejemplo, Google Workspace o Microsoft 365) mediante \textit{tokens}, prescindiendo del almacenamiento de contraseñas.
\end{GlossaryBox}

\begin{GlossaryBox}{RBAC (\textit{Role-Based Access Control})}
    Mecanismo de seguridad y control de acceso basado en roles que restringe o habilita operaciones en el sistema de acuerdo con las responsabilidades y privilegios asignados al perfil del usuario.
\end{GlossaryBox}

\begin{GlossaryBox}{Uptime}
    Métrica porcentual que mide el tiempo de disponibilidad efectiva e ininterrumpida de una infraestructura de software o servicio en la nube durante un periodo determinado (e.g., $99.5\%$ anual o mensual).
\end{GlossaryBox}

\vspace{0.8cm}
\begin{center}
    \small\sffamily\textcolor{MutedText}{--- Fin del Documento de Especificación de Requerimientos ---}
\end{center}
